\documentclass[10pt,a4paper,twocolumn]{article}
\usepackage[utf8]{inputenc}
\usepackage[T5]{fontenc}
\usepackage[vietnamese]{babel}
\usepackage{geometry}
\geometry{left=1.4cm,right=1.4cm,top=1.6cm,bottom=1.6cm,columnsep=0.6cm}
\usepackage{amsmath,amssymb}
\usepackage{booktabs}
\usepackage{array}
\usepackage{xcolor}
\usepackage{graphicx}
\usepackage{tikz}
\usetikzlibrary{shapes.geometric, arrows.meta, positioning, fit, backgrounds, calc}

\definecolor{primaryBlue}{RGB}{21, 101, 192}
\definecolor{accentTeal}{RGB}{0, 121, 107}
\definecolor{accentOrange}{RGB}{230, 81, 0}
\definecolor{accentPurple}{RGB}{106, 27, 154}
\definecolor{softBlue}{RGB}{227, 242, 253}
\definecolor{softGreen}{RGB}{232, 245, 233}
\definecolor{softOrange}{RGB}{255, 243, 224}
\definecolor{softPurple}{RGB}{243, 229, 245}

\title{\Large\textbf{Tổng Hợp Giải Thuật Điều Khiển Mờ Vòng Lồng Kép (Cascaded Dual-FuzzyPD)\\và Hệ Thống Bộ Lọc Xử Lý Tín Hiệu cho Robot Hai Bánh Tự Cân Bằng STM32F407}}
\author{\textbf{Nhóm Phát Triển Hệ Thống Điều Khiển Tự Động \& Robot Tự Hành}\\
\small Vi điều khiển trung tâm: STM32F407VGT6 (ARM Cortex-M4F @ 168\,MHz, Hardware FPU) --- Hệ điều hành thời gian thực: FreeRTOS ($200\,\text{Hz} / 5\,\text{ms}$)}
\date{\small Tháng 11 năm 2026}

\begin{document}

\twocolumn[
\begin{@twocolumnfalse}
\maketitle
\begin{abstract}
\noindent\textbf{\textit{Tóm tắt---}} Tài liệu này tổng hợp toàn bộ hệ thống giải thuật điều khiển, hợp nhất cảm biến và các bộ lọc xử lý tín hiệu đã triển khai trên Robot hai bánh tự cân bằng (\textit{Two-Wheeled Inverted Pendulum --- TWIP}) sử dụng vi điều khiển \textbf{STM32F407VGT6} và \textbf{FreeRTOS} ở tần số đồng bộ $200\,\text{Hz}$ ($T_s = 5\,\text{ms}$). Hệ thống bao gồm: (1) \textbf{Khối tiền xử lý phần cứng \& 4 bộ lọc Kalman không gian trạng thái} (\texttt{kalmanPitch}, \texttt{kalmanYaw}, \texttt{kalmanSpeed\_L/R} có hệ số tắt dần gia tốc $\gamma = 0.80$); (2) \textbf{Giải thuật hợp nhất vận tốc góc Yaw ba chế độ trên miền đạo hàm (Rate-Domain Yaw Fusion)} kết hợp khóa tĩnh ZUPT (\textit{Zero Angular Velocity Update}); (3) \textbf{Chuỗi 5 bộ lọc tách chuyển động Trọng tâm (CoM Decoupling)}, bộ kẹp chống bắn ngược vị trí (\textit{Anti-Slingshot Clamping} $\pm 4.0\,\text{cm}$) và bộ tích phân mặt trượt tự học góc cân bằng tĩnh (\textit{Sliding-Mode Auto-Trim}); (4) \textbf{Kiến trúc điều khiển mờ vòng lồng kép (\texttt{FuzzyPD\_Pos} $\rightarrow$ \texttt{FuzzyPD\_Pitch})} chiếm 90\% công suất PWM ($\pm 7560/8400$) quanh điểm cân bằng chuẩn $\theta_0 = -3.56^\circ$; và (5) \textbf{Bộ điều khiển \texttt{Yaw PID} vi phân kép} (10\% công suất PWM, $\pm 840/8400$), bù vùng chết bất đối xứng 4 góc phần tư có độ dốc chuyển tiếp mượt (\textit{Soft-Knee Deadband}) cùng bộ giới hạn tốc độ biến thiên xung (\textit{Slew-Rate Limiter}). Kết quả thực nghiệm (\texttt{lan39.csv}, $8.063$ mẫu liên tục trong $403\,\text{s}$) chứng minh robot giữ cân bằng tĩnh với độ lệch chuẩn góc $\sigma_\theta = 0.11^\circ$, khóa vị trí trong $\pm 2\,\text{mm}$ và phục hồi hoàn toàn sau các xung lực đẩy mạnh đạt vận tốc đỉnh $78.9\,\text{cm/s}$.\\[4pt]
\noindent\textbf{\textit{Từ khóa---}} Robot hai bánh tự cân bằng (TWIP), Cascaded Dual-FuzzyPD, Bộ lọc Kalman, Hợp nhất cảm biến miền vận tốc góc (Rate-Domain Yaw Fusion), Tách trọng tâm (CoM Decoupling), Chống bắn ngược vị trí (Anti-Slingshot), STM32F407, FreeRTOS.
\end{abstract}
\vspace{6pt}
\hrule
\vspace{8pt}
\end{@twocolumnfalse}
]

\section{Kiến Trúc Đa Nhiệm Thời Gian Thực (FreeRTOS 200 Hz)}
Trên vi điều khiển \textbf{STM32F407VGT6} ($168\,\text{MHz}$, FPU phần cứng), phần mềm nhúng được chia thành 4 tác vụ thời gian thực lập lịch ưu tiên cố định thông qua \texttt{vTaskDelayUntil}:
\begin{itemize}
    \item \texttt{vTaskControl} (Ưu tiên 4 --- cao nhất, $T_s = 5\,\text{ms} / 200\,\text{Hz}$, Stack $1024\,\text{W}$): Thực thi tách chuyển động trọng tâm (CoM), vòng ngoài \texttt{FuzzyPD\_Pos}, tự học góc tĩnh, vòng trong \texttt{FuzzyPD\_Pitch} (90\% PWM), vòng hướng \texttt{Yaw PID} (10\% PWM), bù vùng chết 4 góc phần tư và giới hạn Slew-Rate trước khi nạp vào \texttt{TIM1} ($10\,\text{kHz}$, $\pm 8400\,\text{PWM}$).
    \item \texttt{vTaskIMU} (Ưu tiên 3, $T_s = 5\,\text{ms} / 200\,\text{Hz}$, Stack $512\,\text{W}$): Đọc cảm biến MPU6050 qua \texttt{I2C1} ($400\,\text{kHz}$), tự động hiệu chuẩn điểm 0 tĩnh Gyro Z (ZUPT), hợp nhất vận tốc góc Yaw và cập nhật 2 bộ lọc \texttt{kalmanPitch}, \texttt{kalmanYaw}.
    \item \texttt{vTaskEncoder} (Ưu tiên 3, $T_s = 5\,\text{ms} / 200\,\text{Hz}$, Stack $512\,\text{W}$): Đọc bộ đếm vuông pha $\times 4$ \texttt{TIM3}/\texttt{TIM4} ($1320\,\text{xung/vòng}$), lọc cắt gai xung phần cứng, chạy 2 bộ lọc \texttt{kalmanSpeed\_L/R} và cập nhật Odometry.
    \item \texttt{vTaskUART} (Ưu tiên 2, $T_s = 10\,\text{ms} / 100\,\text{Hz}$, Stack $512\,\text{W}$): Giải mã gói lệnh điều khiển trực tuyến (\texttt{0x01}--\texttt{0x0B}) và truyền khung Telemetry DMA ở tần số $50\,\text{Hz}$.
\end{itemize}

\section{Tiền Xử Lý Phần Cứng \& Các Bộ Lọc Kalman}
\subsection{Bộ lọc phần cứng MPU6050 \& Cắt gai Encoder}
\begin{enumerate}
    \item \textbf{Bộ lọc thông thấp số nội MPU6050 (Hardware DLPF):} Thanh ghi \texttt{CONFIG (0x1A) = 0x03} kích hoạt bộ lọc thông thấp phần cứng với tần số cắt $44\,\text{Hz}$ (Accelerometer $\pm 2g$, $16384\,\text{LSB}/g$) và $42\,\text{Hz}$ (Gyroscope $\pm 500^\circ/\text{s}$, $65.5\,\text{LSB}/(^\circ/\text{s})$), loại bỏ rung cơ khí tần số cao trước khi lấy mẫu. Góc nghiêng tĩnh và vận tốc góc có bảo vệ dấu $\text{sign}(a_z)$:
    \begin{align}
        \theta_{\text{acc}} &= \operatorname{atan2}\!\left(-a_x,\; \sqrt{a_y^2 + a_z^2}\right) \cdot \frac{180^\circ}{\pi} \\
        \dot{\theta}_{\text{gyro}} &= \operatorname{sign}(a_z) \cdot \frac{g_{y,\text{raw}} - g_{y,\text{offset}}}{65.5}
    \end{align}
    \item \textbf{Lọc số đầu vào Timer \& Bộ lọc cắt gai Encoder (Spike Rejection):} Các kênh \texttt{TIM3}/\texttt{TIM4} bật bộ lọc số phần cứng $\texttt{IC1F} = \texttt{IC2F} = \texttt{0x0F}$ trên thanh ghi \texttt{CCMR1}. Sau khi khử tràn 16-bit, mọi xung nhiễu bất thường vượt quá ngưỡng vật lý $|\Delta N(k)| > 45\,\text{xung}/5\,\text{ms}$ đều bị loại bỏ và thay bằng $\Delta N(k-1)$ hợp lệ liền trước:
    \begin{equation}
        v_{\text{raw}}(k) = \frac{60 \cdot \Delta N(k)}{1320 \cdot 0.005} = 9.0909 \cdot \Delta N(k) \quad (\text{RPM})
    \end{equation}
\end{enumerate}

\subsection{Bộ lọc Kalman góc nghiêng 2 trạng thái (\texttt{kalmanPitch}, \texttt{kalmanYaw})}
Vector trạng thái $\hat{\mathbf{x}}_k = [\hat{\theta}_k,\; \hat{b}_k]^T$ ước lượng đồng thời góc thực và độ trôi điểm 0 của Gyroscope với thông số $Q_{\text{angle}} = 0.001$, $Q_{\text{bias}} = 0.003$, $R_{\text{measure}} = 0.03$ và $dt = 0.005\,\text{s}$:
\begin{align}
    \dot{\hat{\theta}}_k &= \omega_{\text{gyro},k} - \hat{b}_{k-1}, \qquad \hat{\theta}_k^- = \hat{\theta}_{k-1} + dt \cdot \dot{\hat{\theta}}_k \\
    \mathbf{K}_k &= \mathbf{P}_k^- \mathbf{H}^T \big(\mathbf{H}\mathbf{P}_k^- \mathbf{H}^T + R_{\text{measure}}\big)^{-1} \\
    \hat{\mathbf{x}}_k &= \hat{\mathbf{x}}_k^- + \mathbf{K}_k \big(\theta_{\text{meas},k} - \hat{\theta}_k^-\big)
\end{align}

\subsection{Bộ lọc Kalman động học vận tốc bánh xe có tắt dần gia tốc (\texttt{kalmanSpeed\_L/R})}
Ở tần số $200\,\text{Hz}$, $1\,\text{xung}$ Encoder tạo ra bước nhảy lượng tử hóa lớn ($9.09\,\text{RPM} \approx 3.09\,\text{cm/s}$). Bộ lọc \texttt{KalmanSpeed\_t} ước lượng vector trạng thái $\hat{\mathbf{x}}_{v,k} = [\hat{v}_k,\; \hat{a}_k]^T$ với \textbf{hệ số tắt dần gia tốc $\gamma = 0.80$} giúp làm phẳng hoàn toàn bậc thang lượng tử hóa với độ trễ pha $< 20\,\text{ms}$ mà không gây vọt lố khi bánh xe đảo chiều ($Q_{\text{speed}} = 0.40$, $Q_{\text{accel}} = 50.0$, $R_{\text{measure}} = 5.0$):
\begin{align}
    \hat{a}_k^- &= 0.80 \cdot \hat{a}_{k-1}, \qquad \hat{v}_k^- = \hat{v}_{k-1} + \hat{a}_k^- \cdot dt \\
    \hat{v}_k &= \hat{v}_k^- + K_{v,0}\big(v_{\text{raw},k} - \hat{v}_k^-\big), \quad \hat{a}_k = \hat{a}_k^- + K_{v,1}\big(v_{\text{raw},k} - \hat{v}_k^-\big)
\end{align}

\section{Hợp Nhất Cảm Biến Hướng Xoay Yaw Trên Miền Đạo Hàm (Rate-Domain Fusion \& ZUPT)}
Nếu trộn trực tiếp góc $\psi_{\text{enc}}$ vào $\psi_{\text{gyro}}$, khi thân xe bị ngoại lực xoay sang hướng mới rồi dừng lại, hai bánh xe đứng yên sẽ giữ $\psi_{\text{enc}}$ ở giá trị cũ và kéo ngược góc Yaw gây sai lệch điều khiển. Do đó, hệ thống thực hiện hợp nhất trên \textbf{miền vận tốc góc (Rate-Domain)} kết hợp \textbf{Cập nhật vận tốc góc bằng 0 (ZUPT)}:
\begin{enumerate}
    \item \textbf{Tự động học điểm 0 tĩnh Gyro Z (Online ZUPT Bias):} Khi bánh xe và thân xe đứng yên ($|\dot{\psi}_{\text{enc}}| < 1.0^\circ/\text{s}$ và $|g_{z,\text{raw}} - g_{z,\text{offset}}| < 80\,\text{LSB}$), điểm 0 được cập nhật qua bộ lọc thông thấp IIR cực chậm:
    \begin{equation}
        g_{z,\text{offset}}(k) = 0.995 \cdot g_{z,\text{offset}}(k-1) + 0.005 \cdot g_{z,\text{raw}}(k)
    \end{equation}
    \item \textbf{Luật hợp nhất vận tốc góc 3 chế độ:}
    \begin{equation}
        \dot{\psi}_{\text{fused}} = \begin{cases}
            0, & |\dot{\psi}_{\text{gyro}}| < 0.7^\circ/\text{s} \land |\dot{\psi}_{\text{enc}}| < 1.2^\circ/\text{s} \\
            0.80\dot{\psi}_{\text{gyro}} + 0.20\dot{\psi}_{\text{enc}}, & \text{khi 2 bánh lăn đồng pha với Gyro Z} \\
            \dot{\psi}_{\text{gyro}}, & \text{khi bị ngoại lực xoay / trượt bánh}
        \end{cases}
    \end{equation}
    Góc $\psi_{\text{raw}}(k) = \psi_{\text{raw}}(k-1) + \dot{\psi}_{\text{fused}} \cdot dt$ sau đó được đưa qua bộ lọc \texttt{kalmanYaw} cùng vận tốc góc $\dot{\psi}_{\text{gyro}}$.
\end{enumerate}

\section{Khử Nhiễu Giao Thoa Trọng Tâm (CoM Decoupling) \& Các Bộ Lọc Vòng Ngoài}
Để triệt tiêu hiện tượng giao thoa không pha cực tiểu (\textit{non-minimum phase coupling}) khi thân xe lắc quanh trục bánh xe, 5 bộ lọc xử lý vòng ngoài được thiết kế nối tiếp:
\begin{enumerate}
    \item \textbf{Tách vận tốc tịnh tiến Trọng tâm (CoM Velocity Decoupling) \& Lọc thông thấp IIR ($\alpha_v = 0.22$):}
    \begin{align}
        v_{\text{CoM,raw}} &= v_{\text{wheel}} - 0.42 \cdot \operatorname{clamp}\big(\dot{\hat{\theta}},\; -60^\circ/\text{s},\; +60^\circ/\text{s}\big) \\
        \hat{v}_{\text{CoM}}(k) &= \hat{v}_{\text{CoM}}(k-1) + 0.22 \cdot \big(v_{\text{CoM,raw}}(k) - \hat{v}_{\text{CoM}}(k-1)\big)
    \end{align}
    \item \textbf{Tách vị trí Trọng tâm (CoM Position Decoupling):}
    \begin{equation}
        x_{\text{CoM}} = x_{\text{fw}} - 0.15 \cdot \big(\hat{\theta} - \theta_{\text{base}}\big) \quad (\text{cm})
    \end{equation}
    \item \textbf{Bộ kẹp chống bắn ngược vị trí (Anti-Slingshot Position Clamping):} Khi đang lái, điểm đích dẫn trước tối đa $\pm 12.0\,\text{cm}$ (\texttt{POS\_TARGET\_LEAD\_MAX\_CM}); ngay khi nhả tay lái, điểm đích chốt tức thời $x_{\text{target}} = x_{\text{CoM}}$ để phanh tại chỗ. Khi đứng yên bị ngoại lực đẩy trượt đi xa, $x_{\text{target}}$ được kẹp trượt theo chân xe trong phạm vi $\pm 4.0\,\text{cm}$ (\texttt{POS\_HOLD\_MAX\_ERR\_CM = 4.0}):
    \begin{equation}
        x_{\text{target}}(k) = \operatorname{clamp}\!\big(x_{\text{target}}(k),\; x_{\text{fw}} - 4.0,\; x_{\text{fw}} + 4.0\big)
    \end{equation}
    Cơ chế này ngăn chặn sai số vị trí tích lũy quá lớn làm nghiêng xe $2^\circ$ bắn ngược về sau (\textit{slingshot}) gây ngã, giúp vòng ngoài ưu tiên tuyệt đối việc hãm vận tốc $v_{\text{CoM}}$ về $0$ tại chỗ.
    \item \textbf{Tự động học góc cân bằng tĩnh theo mặt trượt (Sliding-Mode Static Balance Auto-Trim):} Định nghĩa biến mặt trượt $s_{\text{pos}} = e_{\text{pos}} + 1.2 \cdot v_{\text{CoM,cms}}$ (với $e_{\text{pos}} = x_{\text{fw}} - x_{\text{target}}$). Khâu tích phân chỉ tích lũy khi xe chưa tự hội tụ về điểm cân bằng ($e_{\text{pos}} \cdot s_{\text{pos}} > 0$):
    \begin{equation}
        \theta_{I,\text{pos}}(k) = \operatorname{clamp}\!\big(\theta_{I,\text{pos}}(k-1) - 0.004 \cdot \operatorname{clamp}(s_{\text{pos}}, -4, 4) \cdot dt,\; \pm 1.20^\circ\big)
    \end{equation}
    \item \textbf{Bộ lọc thông thấp IIR góc nghiêng đặt vòng ngoài ($\alpha_{\text{tilt}} = 0.25$):}
    \begin{align}
        \Delta\theta_{\text{filt}}(k) &= \Delta\theta_{\text{filt}}(k-1) + 0.25 \cdot \big(\Delta\theta_{\text{fuzzy}}(k) - \Delta\theta_{\text{filt}}(k-1)\big) \\
        \Delta\theta_{\text{total}}(k) &= \operatorname{clamp}\!\big(\theta_{I,\text{pos}}(k) + \Delta\theta_{\text{filt}}(k),\; -3.20^\circ,\; +3.20^\circ\big)
    \end{align}
\end{enumerate}

\section{Bộ Điều Khiển Mờ Kép Nối Tiếp (Cascaded Dual-FuzzyPD)}
Cả \texttt{FuzzyPD\_Pos} và \texttt{FuzzyPD\_Pitch} đều thực thi qua 5 khâu chuẩn hóa trên $[-1.0, +1.0]$: (1) Tiền xử lý \& kẹp bão hòa $e_{\text{norm}} = \operatorname{sat}(k_e e)$, $de_{\text{norm}} = \operatorname{sat}(k_{de}\dot{e})$; (2) Mờ hóa 5 tập $\{\text{NB}, \text{NS}, \text{ZE}, \text{PS}, \text{PB}\}$ với tập $\text{ZE}$ hình thang cân có vùng phẳng trung tâm $\pm z_e$; (3) Suy diễn tích đại số Larsen $w_{i,j} = \mu_{e,i} \cdot \mu_{de,j}$ trên bảng 25 luật (Bảng~\ref{tab:rules}); (4) Giải mờ trung bình trọng số 7 vạch Singleton $\{-1, -u_m, -u_s, 0, u_s, u_m, 1\}$ và kẹp $[-1, +1]$; (5) Hậu xử lý nhân hệ số khuếch đại $k_u$. Thông số tinh chỉnh tối ưu được tổng hợp tại Bảng~\ref{tab:params}.

\begin{table}[htbp]
\centering
\caption{Ma Trận 25 Luật Mờ $5 \times 5$ Ánh Xạ Ra 7 Mức Đầu Ra Singleton}
\label{tab:rules}
\footnotesize
\begin{tabular}{c|ccccc}
\toprule
$e_{\text{norm}} \backslash de_{\text{norm}}$ & \textbf{NB} & \textbf{NS} & \textbf{ZE} & \textbf{PS} & \textbf{PB} \\
\midrule
\textbf{NB} & NB & NB & NM & NS & ZE \\
\textbf{NS} & NB & NM & NS & ZE & PS \\
\textbf{ZE} & NM & NS & ZE & PS & PM \\
\textbf{PS} & NS & ZE & PS & PM & PB \\
\textbf{PB} & ZE & PS & PM & PB & PB \\
\bottomrule
\end{tabular}
\end{table}

\begin{table}[htbp]
\centering
\caption{Thông Số Bộ Điều Khiển Mờ Kép Trong \texttt{MDK\_ARM/main.c}}
\label{tab:params}
\footnotesize
\begin{tabular}{lcc}
\toprule
\textbf{Tham số cấu hình} & \textbf{\texttt{FuzzyPD\_Pitch} (Vòng trong)} & \textbf{\texttt{FuzzyPD\_Pos} (Vòng ngoài)} \\
\midrule
Chu kỳ lấy mẫu $T_s$ & $5\,\text{ms}$ ($200\,\text{Hz}$) & $5\,\text{ms}$ ($200\,\text{Hz}$) \\
Hệ số chuẩn hóa $k_e$ & $1/11.5\,\text{deg}^{-1}$ ($\pm 11.5^\circ$) & $0.025\,\text{cm}^{-1}$ ($\pm 40.0\,\text{cm}$) \\
Hệ số chuẩn hóa $k_{de}$ & $1/320.0\,(\text{deg/s})^{-1}$ ($\pm 320^\circ/\text{s}$) & $0.060\,(\text{cm/s})^{-1}$ ($\pm 16.7\,\text{cm/s}$) \\
Hệ số khuếch đại $k_u$ & $7560.0\,\text{PWM}$ (90\%) & $2.40^\circ$ \\
Mốc tập mờ $e$: $(e_{\text{nb}}, e_{\text{ns}}, z_e)$ & $(0.46, 0.20, 0.002)$ & $(0.44, 0.16, 0.002)$ \\
Mốc tập mờ $de$: $(de_{\text{nb}}, de_{\text{ns}}, z_{de})$ & $(0.46, 0.20, 0.005)$ & $(0.46, 0.18, 0.010)$ \\
Singleton đầu ra $(u_s, u_m)$ & $(0.25, 0.65)$ & $(0.26, 0.62)$ \\
Độ lợi tuyến tính tương đương & $K_p \approx 822\,\text{PWM}/^\circ, K_d \approx 29.5$ & $K_p \approx 0.098^\circ/\text{cm}, K_d \approx 0.208$ \\
\bottomrule
\end{tabular}
\end{table}

\section{Bộ Điều Khiển Yaw PID Kép \& Xử Lý Công Suất Động Cơ}
\begin{enumerate}
    \item \textbf{Bộ điều khiển \texttt{Yaw PID} vi phân kép (10\% công suất = $\pm 840\,\text{PWM}$):} Kết hợp thành phần cản dịu từ vận tốc vi sai hai bánh xe ($K_{d,\text{wheel}} = 6.30$) và vận tốc góc trực tiếp từ IMU ($K_{d,\text{imu}} = 3.78$), cùng $K_{p,\text{yaw}} = 25.20$, $K_{i,\text{yaw}} = 4.20$ (vùng chết tích phân $\pm 0.35^\circ$, xả nhanh tích phân $I_{\text{yaw}} \leftarrow 0.85 I_{\text{yaw}}$ mỗi khi sai số đảo dấu):
    \begin{equation}
        u_{\text{yaw}} = \operatorname{clamp}\!\big(K_{p,\text{yaw}} e_\psi + K_{i,\text{yaw}} I_\psi + K_{d,\text{wheel}}(-v_L + v_R) + K_{d,\text{imu}}\dot{\hat{\psi}},\; \pm 840\big)
    \end{equation}
    \item \textbf{Bù vùng chết bất đối xứng 4 góc phần tư có độ dốc chuyển tiếp (4-Quadrant Soft-Knee Deadband):} Với ngưỡng ma sát động thực đo ($D_L^+ = 190, D_L^- = 190, D_R^+ = 260, D_R^- = 240\,\text{PWM}$) và vùng chuyển tiếp tuyến tính $U_{\text{knee}} = 120\,\text{PWM}$:
    \begin{equation}
        u_{\text{comp}} = u + \operatorname{sign}(u) \cdot D_{\text{L/R}}^{\pm} \cdot \min\!\left(1.0,\; \frac{|u|}{120.0}\right)
    \end{equation}
    \item \textbf{Giới hạn tốc độ biến thiên xung (PWM Slew-Rate Limiter):} Kẹp bước nhảy $|\Delta\text{PWM}| \le 1000\,\text{PWM}/5\,\text{ms}$ trước khi xuất ra cầu H.
\end{enumerate}

\section{Kết Quả Thực Nghiệm (\texttt{lan39.csv})}
Thử nghiệm thực tế với góc cân bằng chuẩn $\theta_0 = -3.56^\circ$ trong $403\,\text{s}$ liên tục ($8.063$ mẫu Telemetry $50\,\text{Hz}$) ghi nhận: (i) ở trạng thái đứng yên, độ lệch chuẩn góc nghiêng đạt $\sigma_\theta = 0.11^\circ$ (đỉnh-đỉnh $0.51^\circ$), độ lệch chuẩn vận tốc $\sigma_v = 0.70\,\text{cm/s}$, sai số vị trí giữ chặt trong $\pm 0.2\,\text{cm}$ ($\pm 2\,\text{mm}$) và độ trôi góc Yaw bằng $0.000^\circ/\text{s}$; (ii) dưới tác động của hàng loạt xung lực đẩy mạnh liên tiếp từ hai phía (vận tốc bị đẩy đạt đỉnh $78.9\,\text{cm/s}$), robot hoàn toàn không bị bắn ngược vị trí và tự phục hồi cân bằng sau $1.2\text{--}1.8\,\text{s}$ với 0 lần ngã.

\end{document}
